\chapter{Optical surfaces}\label{ch:optics}

Lens and mirror surfaces are specified by their \emph{sag}: the height $z$ of the surface above
the tangent plane at the vertex as a function of the lateral position. \NF{} turns the standard
sag descriptions of optical design -- conics, even and odd aspheres, Zernike polynomials and XY
polynomials, as well as arbitrary sag functions -- into NURBS surfaces that can be trimmed to an
aperture, combined into solids, exported to STEP and meshed for wave-optical or thermal
simulations. Conics are represented exactly; all other sags are fitted to a tolerance that is
verified at dense check points.

\section{Conventions}\label{sec:optics-conventions}

\begin{itemize}
\item The optical axis is $z$, the vertex lies at the origin, and the sag $z(x,y)$ is measured
  along $+z$. Place a surface elsewhere with a \code{Transform} (\cref{ch:transforms}).
\item The curvature is $c = 1/R$. $R > 0$ puts the centre of curvature on the $+z$ side: the
  surface opens towards $+z$ and the sag is positive.
\item The conic constant $k$ selects the conic: $k = 0$ sphere, $k = -1$ paraboloid, $k < -1$
  hyperboloid, $-1 < k < 0$ prolate ellipsoid, $k > 0$ oblate ellipsoid.
\item Lengths are in model units (usually millimetres); all coefficients refer to these units.
\end{itemize}

\section{Rotationally symmetric surfaces}\label{sec:aspheres}

\subsection{The Asphere sag}

\begin{lstlisting}[style=text]
Asphere(curvature=0.0, conic=0.0, even=[], odd=[])
\end{lstlisting}

An \apiref{Asphere} describes the sag
\[
  z(r) = \frac{c\,r^2}{1 + \sqrt{1 - (1+k)\,c^2 r^2}}
         + \sum_{i} A_{2i+4}\, r^{2i+4} + \sum_i B_{i+1}\, r^{i+1},
\]
with \code{even} $= [A_4, A_6, A_8, \dots]$ (coefficients of $r^4, r^6, \dots$) and \code{odd}
$= [B_1, B_2, B_3, \dots]$ (coefficients of $r, r^2, r^3, \dots$). Its members:

\begin{tabularx}{\linewidth}{@{}l>{\raggedright\arraybackslash}X@{}}
\toprule
Member & Meaning \\
\midrule
\code{curvature}, \code{conic}, \code{even}, \code{odd} & the parameters (read/write) \\
\code{sag(r)} & $z(r)$ \\
\code{max_radius()} & the largest $r$ for which the conic term is defined:
  $1/(|c|\sqrt{1+k})$ for $k > -1$, infinite for $k \le -1$ or $c = 0$ \\
\code{is_conic} & \code{True} if there are no polynomial terms (exact representation possible) \\
\bottomrule
\end{tabularx}

\subsection{Profiles and surfaces}

\begin{lstlisting}[style=text]
asphere_profile(asphere, r_max, tolerance=1e-7) -> NurbsCurve
asphere_surface(asphere, r_max, tolerance=1e-7) -> NurbsSurface
\end{lstlisting}

\apiref{asphere_profile} returns the meridian $(r, 0, z(r))$ for $0 \le r \le r_\text{max}$ in the
$xz$ plane; \apiref{asphere_surface} revolves it about the $z$ axis ($u$: azimuth, $v$: radius).

\begin{itemize}
\item \textbf{Pure conics} (\code{is_conic}) are exact: the profile is a piecewise rational
  quadratic (each piece turns the tangent by at most $45^\circ$), and the revolved surface is the
  exact sphere, ellipsoid, paraboloid or hyperboloid -- including hemispheres with a vertical rim
  tangent. Both parameter domains are $[0,1]$.
\item \textbf{Aspheres with polynomial terms} are approximated by a cubic spline profile. The
  profile keeps $x(t) = r$ exactly linear -- the parameter \emph{is} the radius, so the
  $v$ domain of the surface is $[0, r_\text{max}]$ -- and the error is a pure sag error. Knots
  are inserted adaptively until $|\Delta z| \le$ \code{tolerance} at 8 check points per knot
  span; the fit is symmetric, so the slope at the vertex is exactly zero.
\end{itemize}

$r_\text{max}$ must not exceed \code{max_radius()}; beyond it the conic term has no real value.

\example{08_aspheres.py}

Output:

\begin{lstlisting}[style=text]
sphere: is_conic=True, max_radius=40.0, sag(10) = 1.270166537926 (exact 1.270166537926)
  surface: degrees (2, 2), rational, distance from the centre - R: 0.0e+00
conic k = -1.0: max_radius      inf, sag error 1.3e-15
conic k = -0.5: max_radius   28.284, sag error 3.8e-15
conic k =  0.5: max_radius   16.330, sag error 1.8e-15
asphere: is_conic=False, sag(15) = 2.860402219
  profile tol 1e-06: degree 3, 14 ctrl, max distance to the sag 5.8e-07
  profile tol 1e-09: degree 3, 62 ctrl, max distance to the sag 8.6e-10
  surface domains u [0.0, 1.0], v [0.0, 15.0] (v = r for fitted profiles)
  rim point Vec3(12.207391, 8.716629, 2.860402), r = 15.000000000, z - sag(15) = 0.0e+00
  sphere surface domains u [0.0, 1.0], v [0.0, 1.0] (exact conic)
odd asphere sag(10) = 0.100000
\end{lstlisting}

The conic surfaces reproduce the sag to rounding error ($\sim10^{-15}$). For the even asphere a
tolerance of $10^{-6}$ needs 14 control points in the profile, $10^{-9}$ needs 62 -- the number
grows slowly because the cubic spline converges with the fourth power of the knot spacing.

\section{Freeform surfaces}\label{sec:freeform-optics}

\subsection{Zernike polynomials}

\NF{} supports two orderings of the Zernike polynomials $Z_j(\rho, \theta)$, $j \ge 1$, selected
by \apiref{ZernikeIndexing}:

\begin{tabularx}{\linewidth}{@{}l>{\raggedright\arraybackslash}X@{}}
\toprule
Indexing & Definition \\
\midrule
\code{Fringe} & Fringe / University of Arizona ordering, \emph{unnormalised}:
  $Z_4 = 2\rho^2 - 1$ (defocus), $Z_9 = 6\rho^4 - 6\rho^2 + 1$ (primary spherical), $Z_{16}$
  secondary spherical. Common in interferometry. \\
\code{Noll} & Noll ordering, normalised to unit RMS over the unit disc ($Z_4 =
  \sqrt3\,(2\rho^2-1)$). Common in optical design and atmospheric optics. \\
\bottomrule
\end{tabularx}

\apiref{zernike}\code{(j, indexing, rho, theta)} evaluates $Z_j$;
\apiref{zernike_nm}\code{(j, indexing)} returns the radial order $n$ and the azimuthal frequency
$m$ ($m < 0$ denotes a sine term).

\begin{lstlisting}[style=text]
ZernikeSag(base=Asphere(), norm_radius=1.0, indexing=ZernikeIndexing.Fringe, coeffs=[])
zernike_surface(zernike, half_width=0.0, tolerance=1e-7) -> NurbsSurface
\end{lstlisting}

A \apiref{ZernikeSag} adds Zernike terms to a base asphere:
$z(x,y) = z_\text{base}(r) + \sum_j \text{coeffs}[j-1]\, Z_j(r/\text{norm\_radius}, \theta)$.
\code{coeffs[0]} is the coefficient of $Z_1$ (piston), \code{coeffs[3]} that of $Z_4$, and so
on. \code{sag(x, y)} evaluates it. \apiref{zernike_surface} fits the graph $(x, y, z(x,y))$ over
the square $[-h, h]^2$ with $h$ = \code{half_width}, by default \code{norm_radius}.

\subsection{XY polynomials}

\begin{lstlisting}[style=text]
XYPolynomialSag(base=Asphere(), norm_radius=1.0, terms=[])
xy_polynomial_surface(poly, half_width=0.0, tolerance=1e-7) -> NurbsSurface
\end{lstlisting}

An \apiref{XYPolynomialSag} adds $\sum c\,(x/\text{norm})^i (y/\text{norm})^j$ to a base
asphere; \code{terms} is a list of triples \code{(i, j, c)}. \apiref{xy_polynomial_surface}
fits it over a square like \code{zernike_surface}.

\subsection{Arbitrary sag functions}

\begin{lstlisting}[style=text]
sag_surface(f, x0, x1, y0, y1, degree=3, tolerance=1e-7, max_spans=256) -> NurbsSurface
\end{lstlisting}

\apiref{sag_surface} fits the graph of any Python function \code{f(x, y)} over the rectangle
$[x_0,x_1]\times[y_0,y_1]$: diffractive structures, measured deformations given as a function,
or sags from another program. The result is a tensor-product spline of the given degree whose
$x$ and $y$ are exactly linear in $u$ and $v$, so \textbf{$u = x$ and $v = y$}. Knots are
inserted adaptively until $|S_z - f| \le$ \code{tolerance} at $8\times8$ check points per knot
rectangle; polynomials of degree $\le$ \code{degree} are reproduced exactly. \code{GeometryError}
is raised if \code{f} returns a non-finite value (e.g.\ outside the domain of a steep conic) or
more than \code{max_spans} spans per direction would be needed. Zernike and XY surfaces use the
same machinery.

\subsection{Apertures}

Zernike, XY and sag surfaces are graphs over a square. A circular clear aperture is a trimmed
surface with a circular trim loop in the parameter plane -- since $u = x$ and $v = y$ the loop
is simply the aperture circle:

\begin{lstlisting}
aperture = nf.TrimmedSurface(z_surf, nf.circle_loop(0.0, 0.0, 9.0))
\end{lstlisting}

Trimmed surfaces are explained in \cref{ch:trimming}; they are passed exactly to gmsh and to
STEP. For the rotationally symmetric \code{asphere_surface} no trimming is needed -- its rim is
the circle $r = r_\text{max}$.

\subsection{Example}

\example{08_freeform_optics.py}

Output:

\begin{lstlisting}[style=text]
Z1  Fringe (n, m) = (0,  0), Noll (n, m) = (0,  0)
Z4  Fringe (n, m) = (2,  0), Noll (n, m) = (2,  0)
Z9  Fringe (n, m) = (4,  0), Noll (n, m) = (3, -3)
Z11 Fringe (n, m) = (3, -3), Noll (n, m) = (4,  0)
Fringe Z4(rho=1) = 1.0, Noll Z4(rho=1) = 1.732051 (sqrt 3)
zernike surface: domain [-10.0, 10.0] x [-10.0, 10.0], net 35 x 35, degree 3
  max |S_z - sag| at 400 random points: 5.5e-09 (tolerance 1e-8)
xy polynomial: domain [-4.0, 4.0], net 5 x 5, S(2, 1).z = 0.000896000000, sag(2, 1) = 0.000896000000
sag_surface: net 115 x 115, max error 7.6e-07 (tolerance 1e-6)
aperture contains (8, 4): True, (8, 5): False
aperture area 255.290267 (flat disc pi 9^2 = 254.469005)
curvature at the vertex: k1 = 0.01249924, k2 = 0.01249924 (1/80 = 0.01250000)
\end{lstlisting}

The table shows why the indexing must always be stated: $Z_9$ is primary spherical aberration in
Fringe ordering but trefoil in Noll ordering. The XY polynomial with a flat base is a cubic
polynomial and therefore reproduced exactly by the bicubic spline ($5\times5$ control points).
The aperture area is slightly larger than the flat disc because the surface is curved.

The last line shows a limit of fitted surfaces: points are within the tolerance, but derivatives
are less accurate. The curvature at the vertex of the fitted sphere ($R = 80$, default tolerance
$10^{-7}$) is correct to about $6\cdot10^{-5}$ relative. If curvatures matter, fit with a smaller
tolerance -- or use \code{asphere_surface}, which is exact for conics.

\section{Accuracy}\label{sec:optics-accuracy}

\begin{table}[htbp]
\centering\small
\begin{tabularx}{\linewidth}{@{}>{\raggedright\arraybackslash}p{3.6cm}>{\raggedright\arraybackslash}X>{\raggedright\arraybackslash}p{3.8cm}@{}}
\toprule
Surface & Representation & Accuracy \\
\midrule
Sphere, ellipsoid, paraboloid, hyperboloid (\code{Asphere} without terms) & revolved exact
  conic: piecewise rational quadratic profile ($\le 45^\circ$ tangent turn per piece) & exact up
  to rounding (residual $\sim10^{-13}$) \\
Even/odd asphere & cubic spline profile $z(r)$, revolved exactly; $x = r$ exactly linear;
  slope 0 at the vertex & $|\Delta z| \le$ tolerance at 8 check points per span \\
Zernike, XY polynomial, \code{sag_surface} & tensor-product spline graph, $x$, $y$ exactly
  linear; polynomials of degree $\le$ \code{degree} exact & $|\Delta z| \le$ tolerance at
  $8\times8$ check points per knot rectangle \\
\bottomrule
\end{tabularx}
\caption{Representations of optical surfaces.}
\end{table}

The tolerance is \emph{verified at dense check points}, not proven between them. In the test
suite, evaluations at additional random points find errors within 1.5--2 times the tolerance. For
a strict bound, request a smaller tolerance. For steep conic bases, the corners of the square of
a Zernike or XY surface may lie outside the conic's domain (\code{max_radius}); the fit then
reports the undefined sag -- reduce \code{half_width}.

\section{A complete lens}\label{sec:lens}

A lens consists of a front surface, a back surface and the cylindrical edge between their rims.
The following script builds a biconvex lens with an aspheric front and a spherical back surface
and meshes it with gmsh (\cref{ch:meshing}). For a closed solid suitable for Booleans or STEP
export the faces would be sewn into a shell (\cref{ch:solids}).

\example{08_lens_mesh.py}

Output:

\begin{lstlisting}[style=text]
rim heights: front 1.822938, back 5.398492, edge thickness 3.575555
areas: front 462.747309, back 453.526006, edge 269.590472
wrote lens.msh with 1506 nodes
\end{lstlisting}

The back surface has $R = -120$: its centre of curvature lies on the $-z$ side, so after moving
it to $z = 6$ it bulges towards $+z$ and the lens is biconvex. Its rim is lower than its vertex
($5.398 < 6$), the front rim lies at the sag $z(12) = 1.823$, and the edge cylinder connects the
two rims exactly.
